\clearpage
\SourceLeaf{458}
\begin{center}Journal für die reine und angewandte Mathematik herausgegeben von \textsc{Crelle}.\\
Bd. 22. S. 96.\quad Berlin 1841.\end{center}
\begin{center}\textit{Elementare Ableitung eines zuerst von} \textsc{Legendre} \textit{aufgestellten Satzes der sphärischen Trigonometrie.}\end{center}
\noindent\hspace*{1.2em}Sphärische Dreiecke mit kleinen Seiten darf man wie ebene behandeln, wenn man die sphärischen Winkel jeden um den dritten Theil des ganzen sphärischen Excesses vermindert. Die Befugniss dazu lässt sich ganz elementarisch auf folgende Art nachweisen.

Bezeichnet man den ganzen sphärischen Excess eines sphärischen Dreiecks mit $3\omega$; die drei Seiten mit $a,b,c$ und die ihnen gegenüberstehenden sphärischen Winkel mit $A+\omega,B+\omega,C+\omega$; so erhalten ein Paar bekannte Formeln der sphärischen Trigonometrie folgende Gestalt:
\[
\begin{aligned}
\sin\tfrac12 a^2&=\frac{\sin\tfrac32\omega\sin(A-\tfrac12\omega)}{\sin(B+\omega)\sin(C+\omega)}\\[5pt]
\cos\tfrac12 a^2&=\frac{\sin(B-\tfrac12\omega)\sin(C-\tfrac12\omega)}{\sin(B+\omega)\sin(C+\omega)},
\end{aligned}
\]
\noindent aus deren Verbindung folgt
\[
\frac{\sin\tfrac12 a^6}{\cos\tfrac12 a^2}
=\frac{\sin\tfrac32\omega^3\sin(A-\tfrac12\omega)^3}
{\sin(B+\omega)^2\sin(B-\tfrac12\omega)\sin(C+\omega)^2\sin(C-\tfrac12\omega)}.
\]
\noindent Auf gleiche Weise wird
\[
\frac{\sin\tfrac12 b^6}{\cos\tfrac12 b^2}
=\frac{\sin\tfrac32\omega^3\sin(B-\tfrac12\omega)^3}
{\sin(A+\omega)^2\sin(A-\tfrac12\omega)\sin(C+\omega)^2\sin(C-\tfrac12\omega)}.
\]

